% ===================================================================
% Section 5: Computation and tuning
% ===================================================================
\section{Computation and tuning}
\label{sec:computation}

\subsection{Implementation}
The objective \eqref{eq:objective} is assembled from independent \emph{terms},
each exposing a value and a gradient as derived in Lemma~\ref{lem:gradients}; the
covariance is supplied by a swappable estimator, and the constraints of
\eqref{eq:feasible} are composable handlers (a budget equality, per-asset box
bounds, group inequalities, and optional group-target equalities for sector
neutrality). The solver is sequential quadratic programming (SLSQP) with the
analytic Jacobian, warm-started at $\tfrac1N\one$, with a weight floor of
$10^{-12}$ that keeps the barrier terms finite at the boundary and a small
diagonal jitter that keeps $\M^{-1}$ well defined in finite precision.

As set out in Remark~\ref{rem:solver}, global optimality is guaranteed by
convexity rather than by the solver, so we verify the KKT system \eqref{eq:kkt}
at the returned point and check the exact reductions of
Lemma~\ref{lem:reductions} numerically. The reductions hold to solver precision:
the maximum absolute weight deviation from the closed-form GMV portfolio is
$0.0$, and from the ridge-GMV portfolio $3.3\times10^{-8}$. A conic
reformulation---in which the A-term is handled by the Schur-complement
representation of Lemma~\ref{lem:trinv}, the D-term by a log-determinant cone,
and the E-term by the linear matrix inequality $\M\succeq s I$ with $s$
maximized---provides an alternative route that returns a duality-gap certificate
at higher cost.

\subsection{Scaling of the penalties}
\label{sec:penaltyscaling}
Proposition~\ref{prop:scaling} showed that rescaling $V$ is equivalent to
rescaling $\gA$ and $\gE$, which is why Assumption~\ref{as:norm} is imposed.
There is a second, independent scale in the problem: the units of $\Sighat$. The
risk term $\w^\top\Sighat\w$ carries units of squared return, whereas
$\lVert\w\rVert_2^2$ is dimensionless and, under Assumption~\ref{as:norm}, the
three design criteria are dimensionless as well. Comparing a return-squared
quantity with a dimensionless one is meaningless unless one of them is
normalized, so in the implementation the risk and ridge terms are divided by
$\tfrac1N\tr(\Sighat)$, the average asset variance. With this convention the
penalty strengths $(\lambda_2,\gA,\gD,\gE)$ are pure numbers that transfer
across universes of different size, volatility level, and return frequency,
which is what makes a single fixed operating point defensible across the three
studies below.

\subsection{Tuning}
The penalty strengths are hyperparameters. Throughout the experiments we fix a
single interpretable operating point---$\lambda_2=0.10$ and $\gD=\gE=0.05$,
long-only with a $30\%$ per-asset cap---rather than tuning per period, so that
the reported gains cannot be attributed to overfitting the knobs. By
Theorem~\ref{thm:monotone} the qualitative direction of each penalty's effect is
guaranteed regardless of the value chosen, so the operating point governs the
magnitude of the effect and not its sign. A principled alternative is to select
the penalties by cross-validation on a rolling window, minimizing realized
out-of-sample volatility or a turnover-adjusted Sharpe ratio; because the
solution path $\gamma\mapsto\w^\star(\gamma)$ is continuous (and, under strict
complementarity and the conditions of Theorem~\ref{thm:uniqueness},
continuously differentiable by the implicit function theorem applied to
\eqref{eq:kkt}), the path can be traced efficiently by warm-starting along a
grid.

\subsection{Criterion sweeps on synthetic data}
We first verify that each criterion controls the quantity it is designed to
control. The data are generated from a synthetic two-factor return model with
$N=8$ assets and $T=1000$ business days, producing a realistically structured
covariance with condition number $\kappa\approx7.5$. Sweeping one penalty at a
time on the full-sample covariance (long-only) and tracking that criterion's
natural target metric, we find the monotone behaviour predicted by
Theorem~\ref{thm:monotone}: the A-penalty lowers the average estimation variance
$\tr(\M^{-1})$ from $119.9$ to $91.8$; the D-penalty raises the effective number
of assets $\Neff=1/\sum_iw_i^2$ from $6.7$ to $8.0$, i.e.\ to full
diversification; and the E-penalty lowers the condition number $\kappa(\M)$ from
$20.4$ to $14.2$.

\begin{figure}[htbp]
  \centering
  \includegraphics[width=\textwidth]{week6_design_penalty_sweep.png}
  \caption{Criterion sweeps (long-only, synthetic two-factor data). Each design
    penalty improves its own target quantity: A lowers average estimation
    variance, D raises diversification, E improves worst-case conditioning. The
    monotonicity is guaranteed by Theorem~\ref{thm:monotone}.}
  \label{fig:sweep}
\end{figure}

\subsection{High-dimensional simulation ($N/T\to1$)}
\label{sec:highdim}
We next stress the estimation problem directly. On a factor data-generating
process (a low-rank signal plus idiosyncratic noise, the structure that makes the
sample covariance ill-conditioned) we fix the estimation window $T=120$ days and
grow the universe $N\in\{15,40,80,110\}$, so that the concentration ratio
$q=N/T$ rises from $0.125$ to $0.917$. Results are averaged over independent
replications. As $q\to1$ the sample covariance degrades sharply, and the effect
on the portfolios is stark; Table~\ref{tab:nt} reports the extreme case.

\begin{table}[htbp]
\centering
\caption{High-dimensional sweep at $q=N/T=0.917$ ($N=110$, $T=120$), factor
  data-generating process, averaged over replications. Turnover and effective
  $N$ are the stability and diversification readouts.}
\label{tab:nt}
\begin{tabular}{lrrrr}
\toprule
\textbf{Strategy} & \textbf{Sharpe} & \textbf{OOS vol} & \textbf{Turnover} & \textbf{Eff.\ $N$} \\
\midrule
Equal weight              & 4.51 & 0.0347 & 0.000 & 110.0 \\
Markowitz (sample)        & 0.62 & 0.0551 & 3.401 & 8.3 \\
Ridge                     & \textbf{7.35} & \textbf{0.0169} & 0.084 & 101.6 \\
Shrinkage (Ledoit--Wolf)  & 6.25 & 0.0186 & 0.402 & 65.5 \\
Proposed (D$+$E$+$ridge)  & 4.69 & 0.0327 & \textbf{0.0015} & \textbf{110.0} \\
\bottomrule
\end{tabular}
\end{table}

The sample-covariance Markowitz portfolio collapses: turnover $3.40$, fewer than
nine effective names out of $110$, and a Sharpe ratio of $0.62$. Both the
shrinkage and the design-regularized portfolios remain stable. We report the
comparison honestly: as $q\to1$ the design penalty wins decisively on
\emph{turnover and diversification}---turnover $0.0015$ against $0.084$ for
ridge, with near-maximal breadth---while plain ridge attains the best raw
out-of-sample volatility. The two act on different quantities, exactly as the
theory predicts: ridge and the barrier penalties enter the strong-convexity
modulus \eqref{eq:mu} and so damp turnover, whereas only the risk term targets
out-of-sample variance directly.

\begin{figure}[htbp]
  \centering
  \includegraphics[width=\textwidth]{week7_method_vs_nt.png}
  \caption{Behaviour against the concentration ratio $q=N/T$ on a factor
    data-generating process. Left: out-of-sample volatility. Right: average
    turnover. The sample Markowitz portfolio deteriorates as $q\to1$, while the
    structured methods remain stable, with the design-regularized portfolio
    lowest on turnover throughout.}
  \label{fig:nt}
\end{figure}

A companion noise sweep, fixing $N=30$ and $T=120$ and fattening the tails with
Student-$t$ innovations at $\mathrm{df}\in\{3,5,10\}$ against a Gaussian
baseline, gives the same qualitative picture: the structured methods hold their
turnover and breadth as the tails grow heavy, while the raw Markowitz weights
become progressively more erratic. The design-regularized portfolio trades some
risk-adjusted return for this stability under the heaviest tails, an instance of
the trade-off formalized in Theorem~\ref{thm:monotone}.
